% =====================================================================
%  Chương 6 — Học sâu trên đồ thị (làn L6)
% =====================================================================
\chapter{Học sâu trên đồ thị (Deep Graph Learning)}
\label{chap:gnn}

Chương~5 dùng đồ thị tri thức như một \emph{kho} để truy vấn và truy xuất: cạnh được viết ra bởi quy tắc và
được đọc lại bằng phép duyệt. Chương này đi theo hướng ngược lại: đồ thị trở thành \emph{đầu vào của một mô hình học}.
Mạng nơ-ron đồ thị (Graph Neural Network -- GNN) học biểu diễn cho nút, cạnh hoặc cả đồ thị bằng cách để mỗi nút
liên tục trao đổi thông tin với láng giềng. Ta lần lượt đi từ lý do dữ liệu quan hệ cần một kiến trúc riêng, cách biểu
diễn đồ thị bằng ma trận, khung truyền thông điệp tổng quát, ba kiến trúc kinh điển GCN, GraphSAGE, GAT, các giới hạn
khi chồng nhiều lớp, rồi kiểm chứng toàn bộ bằng thực nghiệm trên hai mạng trích dẫn Cora và CiteSeer. Ba lớp GNN đều
được \emph{tự cài} bằng PyTorch thuần (không dùng lớp có sẵn), sau đó đối chiếu với thư viện PyTorch Geometric
\cite{ch6:fey2019pyg} để chắc rằng cài đặt đúng.

Thuật ngữ ``Deep Graph'' ở đây chỉ học sâu trên dữ liệu đồ thị, không đồng nghĩa với Knowledge Graph: KG là một cách
tổ chức tri thức, còn GNN là một họ mô hình học. Một KG có thể được truy vấn bằng SPARQL mà không cần GNN, và một GNN
có thể học trên mạng trích dẫn vốn không được xây như một ontology. Hai hướng gặp nhau ở các bài toán như hoàn thiện
KG (dự đoán cạnh còn thiếu), được trình bày ở mục~\ref{sec:gnn-link}.

% ---------------------------------------------------------------------
\section{Từ dữ liệu dạng lưới đến dữ liệu quan hệ}
\label{sec:gnn-grid}

CNN ở Chương~3 thành công vì ảnh có cấu trúc rất đặc biệt: một lưới đều, mỗi điểm ảnh có đúng tám láng
giềng xếp theo thứ tự cố định trên--dưới--trái--phải. Nhờ thứ tự đó, một kernel $3\times3$ biết trọng số nào nhân với
láng giềng nào và có thể dùng chung trên mọi vị trí. RNN ở Chương~4 cũng dựa vào một thứ tự: trục thời
gian. Rất nhiều dữ liệu thực tế không có thứ tự như vậy: mạng xã hội, mạng trích dẫn khoa học, phân tử hoá học, mạng
giao thông, đồ thị người dùng -- sản phẩm, hay chính đồ thị tri thức. Ở đó, mỗi nút có số láng giềng khác nhau và
không có hệ toạ độ chung (Hình~\ref{fig:ch6-grid}).

\hinhtikz{ch6_grid_vs_graph}{Ảnh là trường hợp riêng của đồ thị: lưới đều với láng giềng có thứ tự. Đồ thị tổng quát
có số láng giềng thay đổi và không có thứ tự toàn cục}{ch6-grid}

Từ quan sát này suy ra yêu cầu cốt lõi với mọi mô hình trên đồ thị. Gọi $P$ là một ma trận hoán vị (đánh số lại các
nút). Đồ thị không thay đổi khi ta đánh số lại, nên:
\begin{itemize}
  \item một hàm dự đoán cho \emph{từng nút} phải \textbf{đồng biến hoán vị} (permutation equivariant):
        $f(PX, PAP^{\top}) = P\,f(X, A)$ -- đánh số lại đầu vào thì đầu ra đánh số lại theo đúng cách đó;
  \item một hàm dự đoán cho \emph{cả đồ thị} phải \textbf{bất biến hoán vị} (permutation invariant):
        $g(PX, PAP^{\top}) = g(X, A)$.
\end{itemize}
Một mạng kết nối đầy đủ nhận vector ghép các đặc trưng nút theo thứ tự chỉ số sẽ vi phạm cả hai tính chất: hoán đổi
nhãn hai nút là đổi đầu ra. Đây là lý do cần các phép tổng hợp không phụ thuộc thứ tự như tổng, trung bình, lấy cực
đại -- nền tảng của khung truyền thông điệp ở mục~\ref{sec:gnn-mp}.

\begin{dinhnghia}{Định nghĩa: Đồ thị có thuộc tính}
Một đồ thị có thuộc tính là bộ $G=(V,E,X)$ với tập nút $V$ ($|V|=N$), tập cạnh $E\subseteq V\times V$ và ma trận
đặc trưng nút $X\in\mathbb R^{N\times d}$ (hàng $x_v$ là đặc trưng của nút $v$). Láng giềng của $v$ là
$\mathcal N(v)=\{u:(u,v)\in E\}$, bậc của $v$ là $d_v=|\mathcal N(v)|$. Cạnh có thể mang thêm đặc trưng $e_{uv}$
(loại quan hệ, trọng số, khoảng cách).
\end{dinhnghia}

% ---------------------------------------------------------------------
\section{Biểu diễn đồ thị bằng ma trận}
\label{sec:gnn-matrix}

Máy tính làm việc với đồ thị qua ba ma trận cơ bản (Hình~\ref{fig:ch6-repr}):
\begin{itemize}
  \item \textbf{Ma trận kề} $A\in\{0,1\}^{N\times N}$, $A_{uv}=1$ nếu có cạnh $u$--$v$. Đồ thị vô hướng cho $A$ đối
        xứng. Thực tế $A$ rất thưa (Cora có hơn hai nghìn bảy trăm nút nhưng mỗi nút chỉ khoảng bốn láng giềng), nên
        được lưu dưới dạng danh sách cạnh \texttt{edge\_index} kích thước $2\times|E|$ hoặc ma trận thưa COO/CSR.
  \item \textbf{Ma trận bậc} $D=\mathrm{diag}(d_1,\dots,d_N)$ với $d_v=\sum_u A_{vu}$.
  \item \textbf{Ma trận Laplacian} $L=D-A$ và dạng chuẩn hoá $L_{\text{sym}}=I-D^{-1/2}AD^{-1/2}$.
\end{itemize}

\hinhtikz{ch6_graph_repr}{Một đồ thị năm nút và ba ma trận biểu diễn nó: ma trận kề, ma trận bậc và Laplacian}{ch6-repr}

Laplacian có ý nghĩa trực quan quan trọng. Với một tín hiệu vô hướng $f\in\mathbb R^N$ gán cho các nút,
\begin{equation}
  f^{\top}Lf=\sum_{(u,v)\in E}(f_u-f_v)^2 ,
  \label{eq:ch6-dirichlet}
\end{equation}
tức là tổng bình phương chênh lệch giữa hai đầu mọi cạnh. Đại lượng này (năng lượng Dirichlet) nhỏ khi tín hiệu
``trơn'' trên đồ thị -- các nút kề nhau có giá trị gần nhau. $L$ đối xứng nửa xác định dương, có phân tích riêng
$L=U\Lambda U^{\top}$; các vector riêng $U$ đóng vai trò ``các sóng Fourier'' của đồ thị, trị riêng nhỏ ứng với tần
số thấp (biến thiên chậm), trị riêng lớn ứng với tần số cao. Phép biến đổi Fourier đồ thị của $f$ là $\hat f=U^{\top}f$.
Đây là cầu nối giữa xử lý tín hiệu và GNN: lọc tín hiệu trên đồ thị nghĩa là nhân phổ $\hat f$ với một hàm $g(\Lambda)$,
\begin{equation}
  g_\theta \star f = U\,g_\theta(\Lambda)\,U^{\top} f .
  \label{eq:ch6-spectral}
\end{equation}
Tính trực tiếp~\eqref{eq:ch6-spectral} cần phân tích riêng $O(N^3)$ và bộ lọc không cục bộ. ChebNet
\cite{ch6:defferrard2016} xấp xỉ $g_\theta$ bằng đa thức Chebyshev bậc $K$ của $L$, khiến bộ lọc chỉ phụ thuộc láng
giềng trong $K$ bước và tính được bằng phép nhân ma trận thưa. GCN ở mục~\ref{sec:gnn-gcn} là trường hợp $K=1$ với
một số đơn giản hoá thêm.

Bảng~\ref{tab:ch6-repr-cost} so sánh các cách lưu đồ thị. Toàn bộ mã của chương dùng \texttt{edge\_index} cho
GraphSAGE và GAT, và ma trận thưa chuẩn hoá $\hat A$ cho GCN.

\begin{table}[H]\centering\caption{Các cách lưu cấu trúc đồ thị và chi phí}\label{tab:ch6-repr-cost}
{\small\begin{tabular}{L{3.6cm}L{2.6cm}L{3.0cm}L{5.0cm}}
\toprule
Cách lưu & Bộ nhớ & Tìm láng giềng & Dùng khi \\
\midrule
Ma trận kề đầy đủ & $O(N^2)$ & $O(N)$ & đồ thị nhỏ, dày; minh hoạ \\
Danh sách cạnh (COO, \texttt{edge\_index}) & $O(|E|)$ & quét $O(|E|)$ & truyền thông điệp bằng \texttt{scatter}/\texttt{index\_add} \\
Ma trận thưa CSR & $O(N+|E|)$ & $O(d_v)$ & nhân $\hat A H$ nhanh trên GPU/CPU \\
Danh sách kề & $O(N+|E|)$ & $O(d_v)$ & lấy mẫu láng giềng (GraphSAGE) \\
\bottomrule
\end{tabular}}
\end{table}

% ---------------------------------------------------------------------
\section{Khung truyền thông điệp}
\label{sec:gnn-mp}

Gilmer và cộng sự \cite{ch6:gilmer2017} chỉ ra rằng hầu hết GNN đều là trường hợp riêng của khung
\textbf{truyền thông điệp} (Message Passing Neural Network -- MPNN). Ở lớp $\ell$, mỗi nút $v$ có biểu diễn $h_v^{(\ell)}$
(khởi đầu $h_v^{(0)}=x_v$) và được cập nhật qua hai bước:
\begin{align}
  m_v^{(\ell)} &= \bigoplus_{u\in\mathcal N(v)} \psi^{(\ell)}\!\big(h_u^{(\ell)}, h_v^{(\ell)}, e_{uv}\big)
  && \text{(AGGREGATE -- tổng hợp thông điệp)} \label{eq:ch6-agg}\\
  h_v^{(\ell+1)} &= \phi^{(\ell)}\!\big(h_v^{(\ell)}, m_v^{(\ell)}\big)
  && \text{(UPDATE -- cập nhật)} \label{eq:ch6-upd}
\end{align}
trong đó $\psi$ tạo thông điệp trên từng cạnh, $\bigoplus$ là phép tổng hợp \emph{không phụ thuộc thứ tự} (tổng,
trung bình, max, hoặc tổng có trọng số attention), $\phi$ là một mạng nhỏ (thường là lớp tuyến tính + phi tuyến).
Vì $\bigoplus$ không quan tâm thứ tự láng giềng và $\psi,\phi$ dùng chung cho mọi nút, mỗi lớp tự động đồng biến hoán
vị. Hình~\ref{fig:ch6-mp} minh hoạ hai lớp truyền chồng lên nhau.

\hinhtikz{ch6_message_passing}{Truyền thông điệp qua hai lớp: ở mỗi lớp, nút đỏ tổng hợp biểu diễn của láng giềng ở
lớp dưới. Sau $L$ lớp, trường tiếp nhận (receptive field) của một nút là toàn bộ láng giềng trong $L$ bước}{ch6-mp}

Ba hệ quả quan trọng rút ra từ khung này:
\begin{enumerate}
  \item \textbf{Trường tiếp nhận tăng theo số lớp}: sau $L$ lớp, $h_v^{(L)}$ phụ thuộc mọi nút cách $v$ không quá $L$
        cạnh. Đây là ``chiều sâu'' của GNN, tương tự việc chồng lớp tích chập mở rộng vùng nhìn của CNN.
  \item \textbf{Chia sẻ tham số}: số tham số không phụ thuộc số nút, nên cùng một mô hình áp dụng được cho đồ thị
        lớn hơn hoặc đồ thị mới -- nền tảng của học \emph{quy nạp} (inductive).
  \item \textbf{Giới hạn biểu đạt}: Xu và cộng sự \cite{ch6:xu2019gin} chứng minh MPNN không mạnh hơn phép thử đẳng cấu
        Weisfeiler--Lehman một chiều; hai nút có cây láng giềng giống nhau sẽ luôn nhận cùng biểu diễn. Tổng hợp bằng
        \emph{tổng} (đơn ánh trên đa tập) đạt đúng giới hạn đó, còn trung bình và max yếu hơn.
\end{enumerate}

Về cài đặt, phép~\eqref{eq:ch6-agg} được hiện thực bằng thao tác \texttt{scatter}: gom thông điệp trên từng cạnh
$(u\to v)$ về chỉ số đích $v$. Trong PyTorch, \texttt{torch.zeros(N,d).index\_add\_(0, dst, msg)} tính tổng thông điệp
của mọi nút trong một lệnh duy nhất, không cần vòng lặp Python -- toàn bộ các lớp tự cài ở mục~\ref{sec:gnn-exp} đều
dựa trên thao tác này.

% ---------------------------------------------------------------------
\section{Ba kiến trúc GNN tiêu biểu}
\label{sec:gnn-arch}

\subsection{GCN -- mạng tích chập đồ thị}
\label{sec:gnn-gcn}

Kipf và Welling \cite{ch6:kipf2017} đi từ bộ lọc phổ~\eqref{eq:ch6-spectral}, giới hạn đa thức Chebyshev ở bậc một và
giả định trị riêng lớn nhất xấp xỉ 2, thu được bộ lọc $g_\theta\star f\approx\theta\,(I+D^{-1/2}AD^{-1/2})f$. Toán tử
$I+D^{-1/2}AD^{-1/2}$ có trị riêng trong $[0,2]$, chồng nhiều lớp dễ làm bùng nổ hoặc triệt tiêu gradient; tác giả dùng
``mẹo chuẩn hoá lại'' (renormalization trick): thêm khuyên vào mọi nút, $\tilde A=A+I$, $\tilde D_{vv}=\sum_u\tilde A_{vu}$,
và định nghĩa lớp GCN cho ma trận đặc trưng nhiều kênh:
\begin{equation}
  H^{(\ell+1)}=\sigma\!\Big(\underbrace{\tilde D^{-1/2}\tilde A\tilde D^{-1/2}}_{\hat A}\,H^{(\ell)}\,W^{(\ell)}\Big).
  \label{eq:ch6-gcn}
\end{equation}
Viết theo từng nút, \eqref{eq:ch6-gcn} chính là một trường hợp của truyền thông điệp:
\begin{equation}
  h_v^{(\ell+1)}=\sigma\Big(\sum_{u\in\mathcal N(v)\cup\{v\}}\frac{1}{\sqrt{\tilde d_u\tilde d_v}}\,W^{(\ell)\top}h_u^{(\ell)}\Big).
\end{equation}
Hệ số $1/\sqrt{\tilde d_u\tilde d_v}$ làm hai việc: nút bậc cao không có biểu diễn lớn bất thường chỉ vì nhiều láng
giềng, và thông điệp từ một ``nút trung tâm'' (hub) có nhiều cạnh bị giảm trọng số. Hình~\ref{fig:ch6-gcn} tách lớp GCN
thành hai bước: \emph{truyền} (nhân $\hat A$) và \emph{biến đổi} (nhân $W$). Thứ tự hai phép nhân có thể đổi; khi
$d'<d$, tính $\hat A(HW)$ rẻ hơn $(\hat AH)W$ -- cài đặt ở Đoạn mã~\ref{code:ch6_gnn__GCNLayer} dùng cách này.

\hinhtikz{ch6_gcn_layer}{Một lớp GCN: trộn đặc trưng theo ma trận kề chuẩn hoá $\hat A$ rồi biến đổi tuyến tính dùng chung}{ch6-gcn}

Nhìn từ góc độ tín hiệu, $\hat A$ là một \emph{bộ lọc thông thấp}: mỗi lần nhân, đặc trưng của nút được thay bằng
trung bình có trọng số với láng giềng, các thành phần tần số cao (chênh lệch giữa nút kề) bị suy giảm. Trên đồ thị
\emph{đồng tính} (homophily cao -- nút kề thường cùng nhãn) như mạng trích dẫn, làm trơn như vậy chính là điều ta
muốn: nhãn của một bài báo được ``bỏ phiếu'' bởi các bài nó trích dẫn. Nhưng lặp lại quá nhiều lần, mọi nút trong một
thành phần liên thông hội tụ về cùng một vector -- hiện tượng over-smoothing ở mục~\ref{sec:gnn-depth}.

\subsection{GraphSAGE -- học quy nạp bằng lấy mẫu và tổng hợp}
\label{sec:gnn-sage}

GCN gốc là \emph{chuyển dẫn} (transductive): $\hat A$ của toàn đồ thị phải có sẵn lúc huấn luyện, và mỗi bước cần đặc
trưng của mọi nút. Hamilton và cộng sự \cite{ch6:hamilton2017} đề xuất GraphSAGE (SAmple and aggreGatE) để học các
\emph{hàm} sinh biểu diễn, áp dụng được cho nút chưa từng thấy:
\begin{align}
  m_v^{(\ell)} &= \mathrm{AGG}^{(\ell)}\big(\{h_u^{(\ell)} : u\in\mathcal S(v)\}\big),\\
  h_v^{(\ell+1)} &= \sigma\big(W_{\text{self}}^{(\ell)}h_v^{(\ell)} + W_{\text{neigh}}^{(\ell)}m_v^{(\ell)}\big),
  \qquad h_v^{(\ell+1)}\leftarrow h_v^{(\ell+1)}/\lVert h_v^{(\ell+1)}\rVert_2 \;(\text{tuỳ chọn}).
\end{align}
Hai khác biệt so với GCN: (i) \textbf{tách nhánh nút gốc} -- $W_{\text{self}}$ riêng cho chính $v$ thay vì trộn $v$
vào trung bình như một láng giềng, giữ lại thông tin riêng của nút; (ii) \textbf{lấy mẫu} $\mathcal S(v)$ -- mỗi lớp
chỉ lấy tối đa $k_\ell$ láng giềng ngẫu nhiên, nên chi phí một nút bị chặn bởi $\prod_\ell k_\ell$ bất kể bậc thật, cho
phép huấn luyện theo mini-batch trên đồ thị hàng trăm triệu nút (PinSage của Pinterest \cite{ch6:ying2018pinsage} là ví
dụ triển khai). Hàm tổng hợp có thể là trung bình (mean), max-pooling sau một lớp MLP, hoặc LSTM trên hoán vị ngẫu
nhiên của láng giềng (Hình~\ref{fig:ch6-sage}).

\hinhtikz{ch6_sage_sampling}{GraphSAGE: lấy mẫu láng giềng có giới hạn ở mỗi bước, tổng hợp từ ngoài vào trong, rồi
kết hợp với nhánh riêng của nút gốc}{ch6-sage}

Cora và CiteSeer đủ nhỏ để dùng toàn bộ láng giềng, nên phần thực nghiệm cài biến thể \emph{mean} không lấy mẫu
(Đoạn mã~\ref{code:ch6_gnn__SAGELayer}); điểm khác GCN còn lại là nhánh $W_{\text{self}}$, làm số tham số tăng gấp đôi.

\subsection{GAT -- trọng số láng giềng học được bằng attention}
\label{sec:gnn-gat}

GCN gán trọng số cố định $1/\sqrt{\tilde d_u\tilde d_v}$, chỉ phụ thuộc bậc. Nhưng không phải láng giềng nào cũng hữu
ích như nhau: một bài báo về mạng nơ-ron có thể trích dẫn cả một bài về lý thuyết xác suất. Veličković và cộng sự
\cite{ch6:velickovic2018} để mô hình tự học trọng số bằng cơ chế attention tương tự Transformer, nhưng chỉ trên láng
giềng thực:
\begin{align}
  e_{vu} &= \mathrm{LeakyReLU}\big(\mathbf a^{\top}[\,W h_v \,\Vert\, W h_u\,]\big), \\
  \alpha_{vu} &= \mathrm{softmax}_{u\in\mathcal N(v)\cup\{v\}}(e_{vu})
               = \frac{\exp(e_{vu})}{\sum_{k\in\mathcal N(v)\cup\{v\}}\exp(e_{vk})},\\
  h_v' &= \big\Vert_{k=1}^{K}\,\sigma\Big(\sum_{u\in\mathcal N(v)\cup\{v\}}\alpha_{vu}^{k}W^{k}h_u\Big).
\end{align}
Chuẩn hoá softmax phải thực hiện \emph{riêng cho từng nút đích} -- nếu dùng một softmax chung cho mọi cạnh, trọng số
của một nút sẽ phụ thuộc vào các cạnh không liên quan. Cài đặt ở Đoạn mã~\ref{code:ch6_gnn__scatter_softmax} trừ giá
trị lớn nhất cục bộ của từng nút trước khi lấy mũ để tránh tràn số. Dùng $K$ đầu attention độc lập rồi ghép lại
(multi-head) giúp ổn định huấn luyện; ở lớp đầu ra, các đầu được lấy trung bình thay vì ghép
(Hình~\ref{fig:ch6-gat}).

\hinhtikz{ch6_gat_attention}{GAT: mỗi cạnh có hệ số $\alpha_{vu}$ học được từ đặc trưng hai đầu; nhiều đầu attention
độc lập được ghép ở lớp ẩn}{ch6-gat}

Cần thận trọng khi diễn giải: trọng số attention cao chỉ cho biết láng giềng đó đóng góp nhiều vào tổng có trọng số,
\emph{không} phải giải thích nhân quả cho quyết định. Mục~\ref{sec:gnn-exp-attn} kiểm tra điều này trên một nút cụ thể.

\subsection{So sánh ba kiến trúc}

\begin{table}[H]\centering\caption{So sánh GCN, GraphSAGE và GAT}\label{tab:ch6-arch}
{\small\begin{tabular}{L{3.0cm}L{3.9cm}L{3.9cm}L{3.9cm}}
\toprule
 & GCN & GraphSAGE & GAT \\
\midrule
Trọng số láng giềng & cố định $1/\sqrt{\tilde d_u\tilde d_v}$ & đều (mean) hoặc học qua pooling/LSTM & học được $\alpha_{vu}$, mỗi đầu một bộ \\
Nút gốc & như một láng giềng (khuyên) & nhánh $W_{\text{self}}$ riêng & khuyên, có attention riêng \\
Chế độ học & chuyển dẫn & quy nạp, mini-batch & quy nạp (cục bộ) \\
Tham số / lớp & $d\,d'$ & $2\,d\,d'$ & $K(d\,d'+2d')$ \\
Chi phí / lớp & $O(|E|d'+Nd\,d')$ & $O(N k\,d')$ khi lấy mẫu $k$ & $O(K(|E|d'+Nd\,d'))$ \\
Điểm mạnh & đơn giản, nhanh, mạnh khi homophily cao & đồ thị rất lớn, nút mới & láng giềng nhiễu, bậc chênh lệch \\
\bottomrule
\end{tabular}}
\end{table}

% ---------------------------------------------------------------------
\section{Các tác vụ học trên đồ thị}
\label{sec:gnn-tasks}

Cùng một bộ mã hoá GNN có thể phục vụ ba mức tác vụ, khác nhau ở chỗ ``đơn vị'' được dự đoán và cách chia tập
(Hình~\ref{fig:ch6-tasks}):
\begin{itemize}
  \item \textbf{Mức nút} -- dự đoán nhãn của từng nút từ biểu diễn $h_v^{(L)}$ qua một lớp softmax. Trong thiết lập
        \emph{bán giám sát chuyển dẫn} của Cora, mô hình thấy toàn bộ cấu trúc và đặc trưng, nhưng hàm mất mát chỉ tính
        trên vài nút có nhãn; nhãn của nút test bị giấu hoàn toàn.
  \item \textbf{Mức cạnh} -- dự đoán xác suất tồn tại (hoặc loại) của cạnh $(u,v)$ từ cặp $(z_u,z_v)$. Phải \emph{xoá}
        cạnh test khỏi đồ thị đầu vào, nếu không mô hình chỉ việc ``nhìn thấy'' đáp án qua truyền thông điệp (rò rỉ).
  \item \textbf{Mức đồ thị} -- dự đoán tính chất của cả đồ thị bằng một phép gộp bất biến hoán vị
        (readout: tổng/trung bình/max của mọi $h_v$). Đồ thị test phải tách khỏi đồ thị train.
\end{itemize}

\hinhtikz{ch6_tasks}{Ba mức tác vụ học trên đồ thị và ví dụ ứng dụng của từng mức}{ch6-tasks}

% ---------------------------------------------------------------------
\section{Over-smoothing, nút thắt và giới hạn độ sâu}
\label{sec:gnn-depth}

Với CNN, chồng thêm lớp (cùng nối tắt như ResNet) gần như luôn có lợi. GNN thì khác: các mô hình tốt nhất trên Cora chỉ
có hai lớp, và thêm lớp thường làm kết quả \emph{giảm}. Li và cộng sự \cite{ch6:li2018} chỉ ra nguyên nhân chính: mỗi
lớp GCN là một bước làm trơn Laplacian; lặp $k$ lần, $\hat A^k X$ hội tụ về một không gian con chỉ phụ thuộc bậc của
nút, mọi nút trong cùng thành phần liên thông trở nên không phân biệt được (Hình~\ref{fig:ch6-smooth}). Hiện tượng này
gọi là \textbf{over-smoothing}.

\hinhtikz{ch6_oversmoothing}{Over-smoothing: sau mỗi lớp truyền, biểu diễn các nút (màu) bị trộn với láng giềng; quá
nhiều lớp làm ba cụm hoà thành một}{ch6-smooth}

Để đo mức làm trơn, Chen và cộng sự \cite{ch6:chen2020mad} đề xuất \textbf{MAD} (Mean Average Distance) -- khoảng cách
cosine trung bình giữa biểu diễn các cặp nút:
\begin{equation}
  \mathrm{MAD}(H)=\frac{1}{N(N-1)}\sum_{u\neq v}\Big(1-\frac{h_u^{\top}h_v}{\lVert h_u\rVert\,\lVert h_v\rVert}\Big).
\end{equation}
MAD tiến về 0 nghĩa là mọi nút đã có cùng hướng biểu diễn. Phần thực nghiệm đo MAD theo số lớp, đồng thời thử một
biện pháp giảm nhẹ là \textbf{nối tắt} (residual) $h^{(\ell+1)}\leftarrow h^{(\ell+1)}+h^{(\ell)}$.

Over-smoothing không phải giới hạn duy nhất:
\begin{itemize}
  \item \textbf{Nút thắt thông tin} (over-squashing) \cite{ch6:alon2021}: số nút trong trường tiếp nhận tăng theo hàm
        mũ của số lớp, nhưng tất cả bị nén vào một vector kích thước cố định; thông tin từ nút ở xa bị ``ép'' qua vài
        cạnh cầu nối.
  \item \textbf{Đồ thị dị tính} (heterophily): khi láng giềng thường \emph{khác} nhãn (mạng hẹn hò, mạng giao dịch gian
        lận), phép làm trơn thông thấp phản tác dụng; cần tách nhánh nút gốc như GraphSAGE hoặc dùng bộ lọc thông cao.
  \item \textbf{Nút bậc cao và đồ thị lớn}: chi phí bộ nhớ của một lớp tỷ lệ với $|E|$; với hàng tỷ cạnh phải lấy mẫu
        hoặc phân vùng đồ thị.
\end{itemize}

\begin{table}[H]\centering\caption{Các giới hạn của GNN sâu và cách kiểm tra}\label{tab:ch6-limits}
{\small\begin{tabular}{L{3.2cm}L{4.8cm}L{6.4cm}}
\toprule
Vấn đề & Hệ quả & Cách đánh giá / giảm nhẹ \\
\midrule
Over-smoothing & biểu diễn nút giống nhau sau nhiều lớp & đo MAD theo số lớp; nối tắt, DropEdge, PairNorm \\
Over-squashing & mất thông tin từ nút ở xa & đo độ nhạy Jacobian; nối lại đồ thị (rewiring) \\
Heterophily & làm trơn trộn lẫn nhãn khác nhau & thống kê homophily; tách nhánh nút gốc \\
Đồ thị thay đổi & biểu diễn cũ thiếu cạnh mới & đánh giá theo thời gian; mô hình quy nạp \\
Nút bậc rất cao & bộ nhớ, độ trễ tăng mạnh & phân bố bậc; lấy mẫu láng giềng \\
\bottomrule
\end{tabular}}
\end{table}

% ---------------------------------------------------------------------
\section{Ứng dụng của học sâu trên đồ thị}
\label{sec:gnn-apps}

\begin{itemize}
  \item \textbf{Đồ thị tri thức}: R-GCN \cite{ch6:schlichtkrull2018} mở rộng GCN với ma trận trọng số riêng cho từng
        loại quan hệ, dùng để phân loại thực thể và hoàn thiện KG (dự đoán bộ ba còn thiếu). Với KG y tế/tài chính ở
        Chương~5, mô hình như vậy có thể đề xuất cạnh ``thuốc -- tương tác -- thuốc'' chưa được ghi,
        nhưng đề xuất đó vẫn phải được kiểm chứng bằng nguồn trước khi đưa vào kho.
  \item \textbf{Hệ gợi ý}: đồ thị hai phía người dùng -- sản phẩm; PinSage \cite{ch6:ying2018pinsage} xử lý đồ thị ba tỷ
        nút bằng lấy mẫu láng giềng theo bước đi ngẫu nhiên.
  \item \textbf{Hoá học và dược học}: phân tử là đồ thị nguyên tử -- liên kết; MPNN \cite{ch6:gilmer2017} dự đoán tính
        chất lượng tử, các mô hình tương tự sàng lọc ứng viên thuốc.
  \item \textbf{Giao thông, mạng xã hội, phát hiện gian lận}: dự báo lưu lượng trên đồ thị đường, phát hiện tài khoản
        bất thường qua mẫu giao dịch.
\end{itemize}
Khảo sát của Wu và cộng sự \cite{ch6:wu2021survey} phân loại đầy đủ các họ GNN (phổ, không gian, tự mã hoá, theo thời
gian); phần còn lại của chương tập trung vào nhánh không gian vì đây là nhánh được dùng rộng rãi nhất.

% =====================================================================
\section{Thực nghiệm: phân loại nút trên mạng trích dẫn}
\label{sec:gnn-exp}

\subsection{Dữ liệu và thiết lập}

Hai bộ dữ liệu chuẩn Planetoid \cite{ch6:sen2008,ch6:yang2016} được dùng: mỗi nút là một bài báo khoa học, đặc trưng là
vector túi từ (bag-of-words) nhị phân trên từ điển của bộ dữ liệu, cạnh là quan hệ trích dẫn (coi như vô hướng), nhãn là
lĩnh vực của bài. Phép chia chuẩn chỉ cho 20 nút có nhãn mỗi lớp để huấn luyện, 500 nút validation và 1.000 nút test --
đây là bài toán \emph{bán giám sát chuyển dẫn}: mô hình thấy toàn bộ đồ thị và đặc trưng, nhưng chỉ thấy nhãn của khoảng
5\,\% số nút (Bảng~\ref{tab:ch6-datasets}).

\bang{ch6_datasets}{Thống kê hai mạng trích dẫn Cora và CiteSeer}{ch6-datasets}

Hai thống kê quyết định GNN có lợi hay không. Thứ nhất, \textbf{homophily cạnh} -- tỷ lệ cạnh nối hai bài cùng lĩnh vực
-- là \SL{ch6_data/Cora/homophilyPct}\,\% trên Cora và \SL{ch6_data/CiteSeer/homophilyPct}\,\% trên CiteSeer: láng giềng
phần lớn cùng nhãn, nên ``hỏi ý kiến láng giềng'' là chiến lược tốt. Thứ hai, đồ thị \textbf{rất thưa}: bậc trung bình
chỉ \SL{ch6_data/Cora/deg_mean} (Cora) và \SL{ch6_data/CiteSeer/deg_mean} (CiteSeer), CiteSeer có
\SL{ch6_data/CiteSeer/isolated} nút cô lập không có cạnh nào -- với những nút này GNN không hơn gì MLP. Phân bố bậc lệch
mạnh (Hình~\ref{fig:ch6-degree}): đa số bài chỉ có 1--4 trích dẫn, vài bài là ``hub'' với hàng chục đến
\SL{ch6_data/Cora/deg_max} cạnh. Các lớp cũng không cân bằng (Hình~\ref{fig:ch6-class}).

\hinhpdf{ch6_class_dist}{1}{Số nút theo lớp: Cora (7 lĩnh vực máy học) và CiteSeer (6 lĩnh vực khoa học máy tính)}{ch6-class}
\hinhpdf{ch6_degree_hist}{1}{Phân bố bậc nút (thang log): đuôi dài, đa số nút có rất ít láng giềng}{ch6-degree}

\subsection{Quy trình và mô hình so sánh}

Hình~\ref{fig:ch6-pipeline} tóm tắt quy trình: nạp đồ thị và chuẩn hoá đặc trưng theo hàng, dựng toán tử đồ thị ($\hat A$
cho GCN, \texttt{edge\_index} cho GraphSAGE/GAT), huấn luyện bằng entropy chéo chỉ trên nút train, chọn epoch có accuracy
validation cao nhất (dừng sớm sau 50 epoch không cải thiện) và báo cáo accuracy test tại epoch đó. Năm mô hình được so
sánh:
\begin{itemize}
  \item \textbf{LogReg}: hồi quy logistic trên vector từ, không dùng cạnh (scikit-learn).
  \item \textbf{MLP}: hai lớp ẩn 64 chiều, dropout 0,5 -- cùng kiến trúc với GCN nhưng \emph{bỏ đồ thị}; đây là phép đối
        chứng trực tiếp cho giá trị của cạnh.
  \item \textbf{GCN} \cite{ch6:kipf2017}: hai lớp, 64 chiều ẩn, dropout 0,5, Adam lr 0,01, weight decay $5\cdot10^{-4}$.
  \item \textbf{GraphSAGE-mean} \cite{ch6:hamilton2017}: cùng cấu hình, tổng hợp trung bình đủ láng giềng.
  \item \textbf{GAT} \cite{ch6:velickovic2018}: lớp 1 có 8 đầu $\times$ 8 chiều (ghép), lớp 2 một đầu, ELU, dropout 0,6
        cho cả đặc trưng và hệ số attention -- đúng cấu hình của bài báo gốc.
\end{itemize}
Mỗi cấu hình chạy với 5 hạt giống (0--4) và báo cáo trung bình $\pm$ độ lệch chuẩn. Ba lớp GNN được tự cài bằng PyTorch
thuần với các phép \texttt{index\_add\_}, \texttt{scatter\_reduce} và nhân ma trận thưa; PyTorch Geometric
\cite{ch6:fey2019pyg} chỉ dùng để tải dữ liệu và để đối chiếu.

\hinhtikz{ch6_pipeline}{Quy trình thực nghiệm phân loại nút: từ đồ thị và đặc trưng đến mô hình, chọn epoch theo
validation và các phân tích sau huấn luyện}{ch6-pipeline}

\maNguon{ch6_gnn__load_planetoid}{Nạp Cora/CiteSeer với phép chia chuẩn và chuẩn hoá đặc trưng theo hàng}
\maNguon{ch6_gnn__gcn_norm_adj}{Ma trận kề chuẩn hoá $\hat A=\tilde D^{-1/2}(A+I)\tilde D^{-1/2}$ dạng thưa}
\maNguon{ch6_gnn__GCNLayer}{Lớp GCN tự cài: biến đổi tuyến tính trước, rồi nhân ma trận thưa $\hat A$}
\maNguon{ch6_gnn__SAGELayer}{Lớp GraphSAGE-mean: nhánh nút gốc và nhánh trung bình láng giềng bằng \texttt{index\_add\_}}
\maNguon{ch6_gnn__scatter_softmax}{Softmax theo từng nút đích trên các cạnh -- phần lõi của GAT}
\maNguon{ch6_gnn__GATLayer}{Lớp GAT nhiều đầu: điểm cạnh, softmax theo láng giềng, tổng có trọng số}

Cài đặt GAT tách vector attention $a=[a_{\text{src}}\,\|\,a_{\text{dst}}]$ thành hai nửa nên điểm của cạnh $(u\to v)$ là
$a_{\text{src}}^{\top}Wh_u + a_{\text{dst}}^{\top}Wh_v$ -- tính một lần cho mọi nút rồi chỉ việc lấy chỉ số theo cạnh, không
phải ghép vector trên từng cạnh. Hàm \texttt{scatter\_softmax} trừ giá trị lớn nhất \emph{của từng nút đích} trước khi lấy
mũ để tránh tràn số. Khuyên tự nối được thêm vào để nút cũng ``chú ý'' tới chính nó.

\maDai{ch6_gnn__GNN}{Mô hình chung: chồng $L$ lớp cùng loại, phi tuyến và dropout giữa các lớp, tuỳ chọn nối tắt}
\maDai{ch6_gnn__train_node}{Huấn luyện bán giám sát: loss trên nút train, chọn epoch theo accuracy validation, dừng sớm}
\maNguon{ch6_run__node_classification}{Vòng thực nghiệm chính: 5 mô hình $\times$ 5 hạt giống $\times$ 2 bộ dữ liệu}
\maNguon{ch6_run__train_one}{Huấn luyện 5 hạt giống cho một cặp (bộ dữ liệu, mô hình) trong tiến trình con riêng}

Việc tách mỗi cặp (bộ dữ liệu, mô hình) ra một tiến trình con là bài học thực tế: trên máy dùng chung với các làn huấn
luyện khác, tiến trình chạy liên tục nhiều giờ từng chết do lỗi truy cập bộ nhớ native ở đúng một điểm (lớp dropout của
GCN trên CiteSeer) trong khi cùng mã đó chạy riêng thì không lỗi. Mỗi cặp xong được lưu thành một tệp, nên một lần chết
chỉ mất phần đang chạy.

\subsection{Kết quả phân loại nút}
\label{sec:gnn-exp-node}

\bang{ch6_node}{Accuracy test (\%, trung bình $\pm$ độ lệch chuẩn qua 5 hạt giống), số tham số và thời gian huấn luyện
trên CPU}{ch6-node}

\hinhpdf{ch6_acc_bars}{1}{Accuracy test của năm mô hình trên hai bộ dữ liệu; thanh lỗi là độ lệch chuẩn qua 5 hạt giống}{ch6-acc}

Bảng~\ref{tab:ch6-node} và Hình~\ref{fig:ch6-acc} trả lời câu hỏi chính: \emph{cạnh có giá trị không?} Trên Cora, cùng
một kiến trúc hai lớp nhưng thêm phép truyền theo đồ thị, GCN đạt \SL{ch6_node/Cora/gcn/test_meanPct}\,\% so với
\SL{ch6_node/Cora/mlp/test_meanPct}\,\% của MLP -- hơn \SL{ch6_summary/Cora/gain_gcn_mlpPct} điểm phần trăm. Trên
CiteSeer, mức tăng là \SL{ch6_summary/CiteSeer/gain_gcn_mlpPct} điểm. Mô hình tốt nhất là \SL{ch6_summary/Cora/best} trên
Cora (\SL{ch6_summary/Cora/best_accPct}\,\%) và \SL{ch6_summary/CiteSeer/best} trên CiteSeer
(\SL{ch6_summary/CiteSeer/best_accPct}\,\%), nhưng khoảng cách giữa GCN và GAT nhỏ, cùng cỡ với độ lệch chuẩn giữa các hạt
giống (Hình~\ref{fig:ch6-seed}). Kết luận thực dụng: phần lớn lợi ích đến từ việc \emph{dùng đồ thị}, không phải từ cách
tổng hợp tinh vi; GraphSAGE, với nhánh nút gốc riêng và gấp đôi tham số, thấp hơn GCN một chút trên cả hai bộ -- trên đồ
thị đồng tính với rất ít nhãn, trộn nút gốc vào trung bình (như GCN) là một dạng điều chuẩn có ích.

Một chi tiết đáng chú ý: trên CiteSeer, hồi quy logistic (\SL{ch6_node/CiteSeer/logreg/test_meanPct}\,\%) \emph{tốt hơn}
MLP (\SL{ch6_node/CiteSeer/mlp/test_meanPct}\,\%). Với 120 nút có nhãn và 3.703 đặc trưng, MLP 64 chiều ẩn quá khớp; mô
hình tuyến tính có điều chuẩn L2 ổn định hơn. GNN tránh được việc này không phải nhờ nhiều tham số hơn (GCN có đúng số
tham số của MLP) mà nhờ \emph{mỗi nút nhận thêm thông tin từ láng giềng}, tương đương với tăng dữ liệu theo cấu trúc.

\hinhpdf{ch6_curves}{1}{Đường học trên Cora (hạt giống 0): loss huấn luyện và accuracy validation theo epoch}{ch6-curves}

Đường học (Hình~\ref{fig:ch6-curves}) cho thấy loss trên 140 nút train của mọi mô hình về gần 0 rất nhanh, nhưng accuracy
validation của MLP dừng ở mức thấp: mô hình đã thuộc lòng tập train mà không tổng quát hoá. GCN và GAT đạt đỉnh validation
sau vài chục epoch rồi đi ngang -- dừng sớm theo validation là cần thiết, vì tiếp tục huấn luyện chỉ làm tăng quá khớp.

\hinhpdf{ch6_seed_box}{1}{Phân bố accuracy test qua 5 hạt giống của từng mô hình (mỗi chấm là một lần chạy)}{ch6-seed}
\hinhpdf{ch6_time_params}{1}{Số tham số và thời gian một lần huấn luyện trên CPU của bốn mô hình học sâu}{ch6-time}

Về chi phí (Hình~\ref{fig:ch6-time}), số tham số của GCN bằng MLP, GraphSAGE gấp đôi (hai ma trận $W_{\text{self}}$,
$W_{\text{neigh}}$), GAT gần bằng GCN dù có 8 đầu vì mỗi đầu chỉ 8 chiều. Thời gian đo trên CPU của một máy đang chạy song
song nhiều tiến trình khác, nên chỉ nên đọc như thứ tự độ lớn: trên đồ thị cỡ vài nghìn nút, một lần huấn luyện GNN mất
cỡ một đến vài phút trên CPU, và GAT đắt hơn GCN vì phải tính hệ số trên từng cạnh của từng đầu.

\hinhpdf{ch6_confusion}{0.78}{Ma trận nhầm lẫn của GCN trên tập test Cora (hạt giống 0)}{ch6-confusion}

Ma trận nhầm lẫn (Hình~\ref{fig:ch6-confusion}, macro-F1 \SL{ch6_extra/confusion/macro_f1}) cho thấy các lỗi tập trung
ở những lớp mang tính ``nền'' và hay được trích dẫn chéo: bài \emph{Lý thuyết} bị nhầm sang \emph{Lập luận theo tình
huống}, bài \emph{Học luật} bị nhầm sang \emph{Lý thuyết}. Các lớp có cộng đồng trích dẫn riêng rõ ràng như mạng nơ-ron
hay thuật toán di truyền được nhận đúng nhiều nhất. Khi cả từ vựng lẫn láng giềng đều pha trộn giữa các lĩnh vực, cả đặc
trưng lẫn cấu trúc đều không đủ để phân biệt.

\subsection{Đối chiếu với thư viện PyTorch Geometric}

Để chắc rằng các lớp tự cài đúng, cùng kiến trúc và siêu tham số được dựng lại bằng \texttt{GCNConv} và \texttt{GATConv}
của PyTorch Geometric (3 hạt giống). Bảng~\ref{tab:ch6-link-pyg} (nửa dưới) cho thấy chênh lệch accuracy trung bình trên
Cora là \SL{ch6_summary/pyg/diff_gcnPct} điểm với GCN và \SL{ch6_summary/pyg/diff_gatPct} điểm với GAT -- nhỏ hơn độ lệch
chuẩn giữa các hạt giống. Phần chênh nhỏ của GCN đến từ chi tiết cài đặt (thứ tự khởi tạo tham số, vòng lặp huấn luyện 300
epoch không dừng sớm của bản PyG), không phải từ công thức.

\maNguon{ch6_run__pyg_crosscheck}{Đối chiếu với GCNConv/GATConv của PyTorch Geometric trên cùng siêu tham số}

% ---------------------------------------------------------------------
\section{Phân tích: độ sâu, cấu trúc và biểu diễn}
\label{sec:gnn-analysis}

\subsection{Over-smoothing khi chồng nhiều lớp}

Mục~\ref{sec:gnn-depth} dự đoán rằng chồng nhiều lớp truyền sẽ làm biểu diễn các nút giống nhau. Thực nghiệm huấn luyện
GCN từ 1 đến 8 lớp (64 chiều ẩn, 3 hạt giống), từ 3 lớp trở lên có thêm phiên bản nối tắt, và đo MAD của biểu diễn ẩn ở
lớp áp chót (Bảng~\ref{tab:ch6-depth}, Hình~\ref{fig:ch6-depth}).

\bang{ch6_depth}{Accuracy test và MAD của GCN theo số lớp, có và không có nối tắt (Cora)}{ch6-depth}
\hinhpdf{ch6_depth}{1}{Accuracy (trục trái) và MAD của biểu diễn ẩn (trục phải) theo số lớp GCN trên Cora; dải mờ là
độ lệch chuẩn qua 3 hạt giống}{ch6-depth}

Hai lớp là tốt nhất (\SL{ch6_summary/depth/best_accPct}\,\%). Một lớp chỉ nhìn láng giềng trực tiếp nên thiếu thông tin;
từ ba lớp trở đi accuracy giảm dần, rồi \emph{sụp đổ} ở 6--8 lớp: GCN 8 lớp không nối tắt chỉ còn
\SL{ch6_summary/depth/acc8Pct}\,\%, thấp hơn cả việc luôn đoán lớp phổ biến nhất, và MAD của nó chỉ còn \SL{ch6_summary/depth/mad8} --
mọi nút gần như có cùng một hướng biểu diễn. Đây đúng là over-smoothing cộng với khó tối ưu của mạng sâu: tín hiệu và
gradient đi qua quá nhiều phép trung bình. Nối tắt giữ lại một phần biểu diễn cũ ở mỗi lớp nên chặn được sự sụp đổ:
ở 8 lớp, bản nối tắt đạt \SL{ch6_summary/depth/acc8_resPct}\,\%, hơn \SL{ch6_summary/depth/res_gain8Pct} điểm. Tuy vậy
nó vẫn không vượt được GCN 2 lớp -- trên Cora, thêm lớp không mang lại thông tin mới đủ để bù chi phí.

\hinhpdf{ch6_propagation}{1}{MAD của $\hat A^k X$ khi lan truyền đặc trưng thô $k$ lần, không có tham số học}{ch6-prop}

Để tách over-smoothing khỏi việc huấn luyện, Hình~\ref{fig:ch6-prop} lặp phép nhân $\hat A$ lên đặc trưng thô mà không có
trọng số nào: MAD giảm đơn điệu từ \SL{ch6_depth/propagation/0/mad} (đặc trưng gốc) xuống \SL{ch6_depth/propagation/7/mad}
sau 64 lần. Đây là bản chất lọc thông thấp của $\hat A$ (mục~\ref{sec:gnn-gcn}): mỗi lần nhân, thành phần ``khác biệt''
giữa nút kề bị san bằng thêm, và giới hạn của $\hat A^k X$ khi $k\to\infty$ chỉ còn phụ thuộc bậc của nút. Mô hình có học
bù được một phần nhờ các ma trận $W$, nhưng không bù được mãi.

\maNguon{ch6_gnn__mad}{Độ đo MAD: khoảng cách cosine trung bình giữa biểu diễn các cặp nút}
\maNguon{ch6_run__depth_study}{Thí nghiệm độ sâu: GCN 1--8 lớp, có/không nối tắt, và lan truyền thuần không học}

\subsection{Giá trị của cấu trúc đồ thị}

GNN thắng MLP vì cạnh mang thông tin, nhưng thông tin đó là gì? Thí nghiệm can thiệp vào cạnh rồi huấn luyện lại GCN:
\emph{bỏ} ngẫu nhiên một tỷ lệ cạnh, hoặc \emph{nối lại} -- giữ nguyên số cạnh nhưng thay một đầu bằng một nút ngẫu nhiên
(Bảng~\ref{tab:ch6-perturb}, Hình~\ref{fig:ch6-perturb}).

\bang{ch6_perturb}{Accuracy GCN trên Cora khi bỏ hoặc nối lại ngẫu nhiên một phần cạnh, và homophily của đồ thị sau can
thiệp}{ch6-perturb}
\hinhpdf{ch6_perturb}{0.95}{Accuracy GCN theo tỷ lệ cạnh bị can thiệp; đường chấm là MLP không dùng cạnh}{ch6-perturb}

Kết quả cho thấy hai điều. Bỏ cạnh làm accuracy giảm từ từ (\SL{ch6_summary/perturb/fullPct}\,\% $\to$
\SL{ch6_summary/perturb/drop50Pct}\,\% khi bỏ một nửa) và khi bỏ hết thì GCN chỉ còn
\SL{ch6_summary/perturb/drop100Pct}\,\% -- trở về mức của MLP, đúng như dự đoán vì $\hat A$ khi đó chỉ còn khuyên tự nối.
Nhưng \emph{cạnh sai còn tệ hơn không có cạnh}: nối lại toàn bộ cạnh ngẫu nhiên kéo homophily xuống
\SL{ch6_summary/perturb/homo_shuf100Pct}\,\% và accuracy xuống \SL{ch6_summary/perturb/shuf100Pct}\,\%, thấp hơn hẳn
MLP (chênh lệch \SL{ch6_summary/perturb/shuf100_vs_mlpPct} điểm). Mô hình vẫn ``tin'' láng giềng và trung bình hoá đặc trưng của các bài
báo không liên quan, làm hỏng chính đặc trưng của nút. Bài học thực tế: trước khi dùng GNN, hãy kiểm tra homophily và chất
lượng cạnh; một đồ thị nhiễu có thể làm mô hình kém hơn việc bỏ qua đồ thị.

\maNguon{ch6_gnn__homophily}{Homophily cạnh: tỷ lệ cạnh nối hai nút cùng nhãn}
\maNguon{ch6_run__edge_perturbation}{Can thiệp vào cấu trúc: bỏ cạnh hoặc nối lại đầu cạnh ngẫu nhiên, huấn luyện lại GCN}

\subsection{Nút nào được lợi từ đồ thị?}

\hinhpdf{ch6_degree_acc}{1}{Accuracy test trên Cora theo nhóm bậc của nút: MLP, GCN và GAT}{ch6-degree-acc}

Chia nút test theo bậc (Hình~\ref{fig:ch6-degree-acc}): MLP gần như không đổi theo bậc vì nó không nhìn cạnh, còn GCN và
GAT tốt dần khi nút có nhiều láng giềng. Với nút chỉ có một láng giềng, GCN hơn MLP \SL{ch6_summary/degree/low_gainPct}
điểm; với nút từ 9 láng giềng trở lên, mức hơn là \SL{ch6_summary/degree/high_gainPct} điểm. Nhiều láng giềng nghĩa là
nhiều ``phiếu bầu'' hơn, và trên đồ thị đồng tính, đa số phiếu thường đúng. Ngược lại, đây cũng là điểm yếu: nút mới, ít
liên kết (bài báo vừa xuất bản, người dùng mới) nhận được ít lợi ích nhất -- bài toán ``khởi đầu lạnh'' quen thuộc của hệ
gợi ý.

\subsection{Không gian biểu diễn}

\hinhpdf{ch6_tsne}{1}{t-SNE của đặc trưng thô (trái), biểu diễn ẩn của GCN (giữa) và của GAT (phải) trên Cora, tô màu
theo nhãn thật}{ch6-tsne}

Hình~\ref{fig:ch6-tsne} chiếu ba không gian về hai chiều bằng t-SNE \cite{ch6:maaten2008}. Đặc trưng túi từ thô gần như
không tách lớp (hệ số silhouette theo cosine \SL{ch6_extra/tsne/raw/silhouette}); sau một lớp GCN, các lớp thành cụm rõ
(silhouette \SL{ch6_extra/tsne/gcn/silhouette}), GAT tương tự (\SL{ch6_extra/tsne/gat/silhouette}). Lưu ý rằng mô hình
chỉ thấy nhãn của 140 nút: cụm hình thành được là nhờ lan truyền đặc trưng qua cạnh, không phải nhờ giám sát dày đặc.

\subsection{GAT thực sự ``chú ý'' vào đâu?}
\label{sec:gnn-exp-attn}

\hinhpdf{ch6_gat_attn}{1}{Hệ số attention lớp 1 (trung bình 8 đầu; chấm đen là từng đầu) của GAT cho các láng giềng của
một nút test bậc vừa trên Cora}{ch6-gat-attn}

Mục~\ref{sec:gnn-gat} cảnh báo rằng hệ số attention không nhất thiết là lời giải thích. Hình~\ref{fig:ch6-gat-attn} lấy
nút test \SL{ch6_extra/gat_attention/node} có \SL{ch6_summary/gat/n_neighbors} láng giềng (kể cả chính nó). Kết quả bất
ngờ nhưng có tính đại diện: các hệ số gần như \emph{bằng nhau} (lớn nhất chỉ \SL{ch6_summary/gat/max_alpha}, xấp xỉ
$1/8$), và cả tám đầu cho cùng một phân bố. Trên Cora, GAT học được attention gần đều -- tức là gần với một phép trung bình
như GCN -- và điều đó giải thích vì sao hai mô hình có accuracy gần như nhau. Nhận xét này trùng với các phân tích sau
bài báo gốc: trên đồ thị đồng tính, ít nhiễu, chẳng có láng giềng nào cần bị ``tắt'', nên attention không có gì để học.
Nó cũng cho thấy không nên đọc một bản đồ attention như lời giải thích mà không kiểm tra xem nó có khác đều hay không.

\maNguon{ch6_run__gat_attention}{Trích hệ số attention của GAT cho các láng giềng của một nút}

% ---------------------------------------------------------------------
\section{Dự đoán liên kết bằng Graph Auto-Encoder}
\label{sec:gnn-link}

Bài toán mức cạnh hỏi: \emph{hai nút chưa nối có nên được nối không?} -- chính là hoàn thiện đồ thị tri thức hay gợi ý
trích dẫn. Graph Auto-Encoder \cite{ch6:kipf2016gae} dùng một GCN hai lớp làm bộ mã hoá ra vector $z_v$ cho mỗi nút và bộ
giải mã đơn giản $\hat p(u,v)=\sigma(z_u^{\top}z_v)$ (Hình~\ref{fig:ch6-gae}). Cạnh của Cora được chia 85/5/10\,\% thành
train/validation/test (\SL{ch6_link/train_edges} cạnh train, \SL{ch6_link/test_edges} cạnh test); bộ mã hoá chỉ thấy cạnh
train, mỗi epoch lấy ngẫu nhiên số mẫu âm (cặp không nối) bằng số cạnh dương, và loss là entropy chéo nhị phân.

\hinhtikz{ch6_gae}{Graph Auto-Encoder: bộ mã hoá GCN tạo biểu diễn nút, bộ giải mã tích vô hướng chấm điểm từng cặp
nút}{ch6-gae}

\maNguon{ch6_gnn__GAE}{Graph Auto-Encoder: bộ mã hoá GCN hai lớp và bộ giải mã tích vô hướng}
\maDai{ch6_run__link_prediction}{Huấn luyện GAE: chia cạnh, lấy mẫu âm mỗi epoch, chọn mô hình theo AUC validation}

\bang{ch6_link_pyg}{Dự đoán liên kết bằng GAE (3 hạt giống) và đối chiếu lớp tự cài với PyTorch Geometric}{ch6-link-pyg}
\hinhpdf{ch6_link_curve}{1}{Loss huấn luyện và AUC validation của GAE theo epoch (hạt giống 0)}{ch6-link}

GAE đạt AUC test \SL{ch6_link/test_auc_meanPct}\,\% và AP \SL{ch6_link/test_ap_meanPct}\,\%
(Bảng~\ref{tab:ch6-link-pyg}): với một cặp (cạnh thật, cặp ngẫu nhiên), mô hình xếp cạnh thật cao hơn trong khoảng chín
trên mười trường hợp, dù không có nhãn chủ đề nào được dùng. AUC validation tăng nhanh trong khoảng 100 epoch đầu rồi tăng chậm dần
(Hình~\ref{fig:ch6-link}). Cần đọc con số này đúng bối cảnh: mẫu âm là cặp ngẫu nhiên nên phần lớn rất dễ loại (hai bài
khác hẳn chủ đề); trong ứng dụng thật, như hoàn thiện KG ở Chương~5, mẫu âm khó hơn nhiều và mọi cạnh được đề xuất vẫn
phải được kiểm chứng bằng nguồn.

% ---------------------------------------------------------------------
\section{Triển khai mô hình lên Web App}
\label{sec:gnn-deploy}

Mô hình GCN, GraphSAGE, GAT và MLP huấn luyện trên Cora (hạt giống 0) được lưu và phục vụ qua \texttt{tl3/serve\_ch6.py}.
Trải nghiệm \texttt{ch6\_cora} có kiểu ``graph'': người dùng nhập mã một bài báo (nút 0--2707) và chọn mô hình; hàm
\texttt{predict} trả về chủ đề dự đoán, phân bố xác suất trên 7 chủ đề, nhãn thật, nút thuộc tập nào (train, validation,
test hay không gán nhãn), bậc của nút và danh sách láng giềng kèm nhãn của chúng. Vì GNN chuyển dẫn tính biểu diễn cho cả
đồ thị một lượt, mô-đun chạy mô hình trên toàn bộ Cora một lần khi nạp (trên CPU, đồ thị nhỏ) rồi chỉ tra kết quả cho từng
yêu cầu. Hiển thị láng giềng cùng nhãn giúp người dùng thấy trực tiếp cơ chế ``hỏi ý kiến láng giềng'': khi đa số láng
giềng cùng chủ đề, GNN đúng ngay cả với nút mà MLP sai.

\maNguon{serve_ch6__SPECS}{Khai báo trải nghiệm \texttt{ch6\_cora} cho Web App}
\maDai{serve_ch6__predict}{Hàm suy luận: mã nút $\to$ chủ đề, xác suất, láng giềng và nhãn thật}

\anhUI{ch6_cora_web}{Phân loại bài báo trên đồ thị trích dẫn Cora trong Web App (desktop): chủ đề dự đoán, phân bố xác
suất và láng giềng}{ch6-ui-web}
\anhMobile{ch6_cora_mobile}{Trải nghiệm phân loại nút Cora trên khung điện thoại 390 px}{ch6-ui-mobile}

% ---------------------------------------------------------------------
\section{Kết luận chương}

Chương này đi từ lý do dữ liệu quan hệ cần một kiến trúc riêng, qua ba ma trận biểu diễn đồ thị và khung truyền thông
điệp, tới ba kiến trúc GCN, GraphSAGE, GAT -- tự cài bằng PyTorch thuần và khớp với thư viện chuẩn. Thực nghiệm trên hai
mạng trích dẫn cho các kết luận có số liệu:
\begin{itemize}
  \item \textbf{Cạnh có giá trị lớn}: cùng kiến trúc, dùng đồ thị nâng accuracy thêm \SL{ch6_summary/Cora/gain_gcn_mlpPct}
        điểm trên Cora và \SL{ch6_summary/CiteSeer/gain_gcn_mlpPct} điểm trên CiteSeer khi chỉ có 20 nhãn mỗi lớp.
  \item \textbf{Kiến trúc ít quan trọng hơn dữ liệu}: GAT tốt nhất nhưng chỉ hơn GCN trong phạm vi độ lệch chuẩn; attention
        học được gần như đều.
  \item \textbf{Sâu hơn không tốt hơn}: hai lớp là tối ưu; tám lớp không nối tắt sụp về \SL{ch6_summary/depth/acc8Pct}\,\%
        với MAD gần 0, nối tắt cứu lại được phần lớn.
  \item \textbf{Cạnh sai còn tệ hơn không có cạnh}: nối lại ngẫu nhiên toàn bộ cạnh khiến GCN kém MLP; homophily là điều
        kiện cần kiểm tra trước khi chọn GNN.
  \item \textbf{Dự đoán liên kết} không cần nhãn: GAE đạt AUC \SL{ch6_link/test_auc_meanPct}\,\%, nền tảng cho hoàn thiện
        đồ thị tri thức.
\end{itemize}
Các giới hạn: thực nghiệm dùng hai đồ thị nhỏ, đồng tính, chuyển dẫn; trên đồ thị dị tính, đồ thị rất lớn hoặc thay đổi
theo thời gian, thứ tự các mô hình có thể khác và cần lấy mẫu láng giềng như GraphSAGE.
